% sbstyle.tex: one style for every figure of the paper, diagrams and data figures alike (proposal).
% Fonts are the paper's own fonts under pdfLaTeX (acmart sigconf): Linux Libertine/Biolinum, Inconsolata, newtxmath,
% so figure text matches the body text. Every figure is built at its printed size and included at 1:1.
\usepackage[T1]{fontenc}
\usepackage[tt=false,type1=true]{libertine}
\usepackage[varqu]{zi4}
\usepackage[libertine]{newtxmath}
\usepackage{pifont}
\usepackage{textcomp}
\usepackage{booktabs}
\usepackage{colortbl}
\usepackage{tikz}
\usepackage{pgfplots}
\usepackage{pgfplotstable}
\pgfplotsset{compat=1.18}
\usetikzlibrary{arrows.meta,calc,positioning,fit,backgrounds,shapes.geometric,patterns}
\usepgfplotslibrary{groupplots}

% ---------------------------------------------------------------- sizes
\newlength{\sbcol}\setlength{\sbcol}{241.14749pt}   % \columnwidth of acmart sigconf
\newlength{\sbwide}\setlength{\sbwide}{506.295pt}   % \textwidth
% One text size, 7.5 pt: Biolinum's x-height is 17% below Arial's, so 7.5 pt Biolinum reads close to today's 6.5-7 pt Arial
% while the paper keeps its page budget (no figure taller than today). Hierarchy by weight, not size. Code and counts inside
% bars at 7 pt (Inconsolata 7 pt has the x-height of Biolinum 7.5 pt).
\newcommand{\sbtext}{\sffamily\fontsize{7.5}{9}\selectfont}
\newcommand{\sbnum}{\sffamily\fontsize{7}{8.4}\selectfont}
\newcommand{\sbcode}{\ttfamily\fontsize{7}{8.4}\selectfont}
\newcommand{\sbtt}[1]{{\ttfamily\fontsize{7}{9}\selectfont #1}}% code inside running figure text
\newcommand{\sbtitle}[1]{{\sbtext\bfseries #1}}
\newcommand{\sbboxtext}{\sffamily\fontsize{7.5}{8.7}\selectfont}% text inside diagram boxes

% ---------------------------------------------------------------- colors
% Rule: a figure uses ONE family -- models (when the series are models), or ink/gray/red (when the series are outcomes,
% methods or stages). Blue/orange mean interpretations A/B and nothing else.
% neutrals: text, axes, and the wanted outcome or our method (ink); baselines and other categories (grays)
\definecolor{sbInk}{HTML}{262626}
\definecolor{sbDark}{HTML}{4A4A4A}
\definecolor{sbMid}{HTML}{8C8C8C}
\definecolor{sbLight}{HTML}{C8C8C8}
\definecolor{sbFaint}{HTML}{E4E4E4}
\definecolor{sbPale}{HTML}{F1F1F1}
\definecolor{sbGrid}{HTML}{E2E2E2}
% the one alert color: failure, a silent choice, a wrong alternative
\definecolor{sbBad}{HTML}{C8432F}
\definecolor{sbBadL}{HTML}{EDB2A6}
% interpretations A and B (diagrams, and wherever the two interpretations themselves are shown)
\definecolor{sbA}{HTML}{2F6EB5}\definecolor{sbAf}{HTML}{DCE8F6}
\definecolor{sbB}{HTML}{D9822B}\definecolor{sbBf}{HTML}{FBE6CF}
% models: hue = vendor (Claude purple, GPT green), lightness = tier (lighter = smaller model)
\definecolor{sbHaiku}{HTML}{9A7BC6}
\definecolor{sbSonnet}{HTML}{55318F}
\definecolor{sbSol}{HTML}{176B43}
\definecolor{sbLuna}{HTML}{52A267}

% ---------------------------------------------------------------- legend keys (one look in every figure)
\newcommand{\sbkeydot}[2]{\tikz[baseline=-0.55ex]\fill[#1] (0,0) circle (1.7pt);\hspace{2pt}#2}
\newcommand{\sbkeyring}[2]{\tikz[baseline=-0.55ex]\draw[#1, line width=0.6pt] (0,0) circle (1.45pt);\hspace{2pt}#2}
\newcommand{\sbkeybox}[2]{\tikz[baseline=-0.55ex]\path[#1] (0,-2.6pt) rectangle (6.5pt,2.6pt);\hspace{2.5pt}#2}
\newcommand{\sbkeymark}[2]{\tikz[baseline=-0.55ex]\path (0,0) pic{sbmark={#1}};\hspace{2pt}#2}
\newcommand{\sbkeydiamond}[2]{\tikz[baseline=-0.55ex]\fill[#1] (0,2.3pt) -- (1.8pt,0) -- (0,-2.3pt) -- (-1.8pt,0) -- cycle;\hspace{2pt}#2}
\newcommand{\sbkeyopendiamond}[2]{\tikz[baseline=-0.55ex]\draw[#1, line width=0.6pt] (0,2.3pt) -- (1.8pt,0) -- (0,-2.3pt) -- (-1.8pt,0) -- cycle;\hspace{2pt}#2}
\newcommand{\sbkeysquare}[2]{\tikz[baseline=-0.55ex]\fill[#1] (-1.6pt,-1.6pt) rectangle (1.6pt,1.6pt);\hspace{2pt}#2}
\newcommand{\sbkeytriangle}[2]{\tikz[baseline=-0.55ex]\fill[#1] (90:2.2pt) -- (210:2.2pt) -- (330:2.2pt) -- cycle;\hspace{2pt}#2}
\newcommand{\sbkeystar}[2]{\tikz[baseline=-0.55ex]\path plot[mark=sbstar, mark size=2.6pt, mark options={fill=#1, draw=#1}] coordinates{(0,0)};\hspace{2pt}#2}
\newcommand{\sbsep}{\hspace{8pt}}
\newcommand{\sbmodelkeys}{\sbkeydot{sbHaiku}{Haiku}\sbsep\sbkeydot{sbSonnet}{Sonnet}\sbsep\sbkeydot{sbSol}{GPT-Sol}\sbsep\sbkeydot{sbLuna}{GPT-Luna}}
\tikzset{sblegend/.style={font=\sbtext, inner sep=0pt, align=left, anchor=south west},
         pics/sbmark/.style={code={\path plot[mark=#1, mark size=2pt] coordinates{(0,0)};}}}
\newcommand{\yes}{{\color{sbInk}\ding{51}}}\newcommand{\no}{{\color{sbBad}\ding{55}}}

% ---------------------------------------------------------------- data figures
\pgfplotsset{
  sb/.style={
    scale only axis, clip=false,
    font=\sbtext, tick label style={font=\sbtext, text=sbInk, execute at begin node={\hyphenpenalty=10000\exhyphenpenalty=10000\relax}}, label style={font=\sbtext, text=sbInk},
    title style={font=\sbtext\bfseries, anchor=south west, at={(0,1)}, inner sep=0pt, yshift=4pt},
    axis lines*=left, axis line style={draw=sbInk, line width=0.5pt},
    every tick/.style={draw=sbInk, line width=0.5pt}, major tick length=2pt, tick align=outside,
    grid style={draw=sbGrid, line width=0.4pt},
    xlabel style={yshift=2pt}, ylabel style={yshift=-3pt},
    every axis plot/.append style={line width=0.6pt},
    legend style={draw=none, fill=none, font=\sbtext},
    % numbers in the text font (Biolinum), not in math digits
    /pgf/number format/assume math mode=true, /pgf/number format/1000 sep={,},
    xticklabel={\pgfmathprintnumber{\tick}}, yticklabel={\pgfmathprintnumber{\tick}},
  },
  % categories on y: no y axis line or ticks, a light grid on the value axis
  sb ycats/.style={y axis line style={draw=none}, ytick style={draw=none}, xmajorgrids, yticklabel style={align=right}},
  sb xcats/.style={x axis line style={draw=sbInk}, xtick style={draw=none}, ymajorgrids},
  % intervals: no caps, the color of their series
  sb xerr/.style={error bars/x dir=both, error bars/x explicit, error bars/error mark=none, error bars/error bar style={line width=0.6pt}},
  sb yerr/.style={error bars/y dir=both, error bars/y explicit, error bars/error mark=none, error bars/error bar style={line width=0.5pt}},
  % points: filled, with a thin white rim so overlapping points stay readable
  sb pt/.style={only marks, mark=*, mark size=1.7pt, color=#1, mark options={fill=#1, draw=white, line width=0.25pt}},
  sb ring/.style={only marks, mark=o, mark size=1.45pt, color=#1, mark options={draw=#1, fill=white, line width=0.6pt}},
  % bars: flat fills, no outlines; stacked segments are split by a hairline of white
  sb bar/.style={fill=#1, draw=white, line width=0.35pt},
}
% four models as points with intervals, one row per y; CSV columns h,hm,hp, s,sm,sp, g,gm,gp, l,lm,lp
\newcommand{\sbmodeldots}[2][0.075]{%
  \addplot[sb pt=sbHaiku,  sb xerr] table[x=h, y expr=\thisrow{y}+3*#1, x error minus=hm, x error plus=hp] {#2};
  \addplot[sb pt=sbSonnet, sb xerr] table[x=s, y expr=\thisrow{y}+#1,   x error minus=sm, x error plus=sp] {#2};
  \addplot[sb pt=sbSol,    sb xerr] table[x=g, y expr=\thisrow{y}-#1,   x error minus=gm, x error plus=gp] {#2};
  \addplot[sb pt=sbLuna,   sb xerr] table[x=l, y expr=\thisrow{y}-3*#1, x error minus=lm, x error plus=lp] {#2};}
% counts centred in the segments of stacked bars, drawn on an overlay axis with the frame of the named axis; CSV columns y,xc,lab,tc
\newcommand{\sbseglabels}[3]{%
  \begin{axis}[sb, hide axis, at={(#1.south west)}, anchor=south west, #2]
  \addplot[only marks, mark=none, point meta=explicit symbolic, visualization depends on={value \thisrow{tc} \as \tc},
    nodes near coords, nodes near coords style={font=\sbnum, text=\tc, anchor=center, inner sep=0pt}]
    table[x=xc, y=y, meta=lab, col sep=comma] {#3};
  \end{axis}}
\pgfplotsset{table/col sep=comma,
  % keep only the rows whose column #1 equals #2 (string compare, so it works for any CSV column)
  discard if not/.style 2 args={x filter/.append code={\edef\sbtempa{\thisrow{#1}}\edef\sbtempb{#2}\ifx\sbtempa\sbtempb\else\def\pgfmathresult{nan}\fi}},
  discard if/.style 2 args={x filter/.append code={\edef\sbtempa{\thisrow{#1}}\edef\sbtempb{#2}\ifx\sbtempa\sbtempb\def\pgfmathresult{nan}\fi}}}
\pgfdeclareplotmark{sbstar}{%
  \pgfpathmoveto{\pgfqpointpolar{90}{1.35\pgfplotmarksize}}%
  \pgfpathlineto{\pgfqpointpolar{126}{0.55\pgfplotmarksize}}%
  \pgfpathlineto{\pgfqpointpolar{162}{1.35\pgfplotmarksize}}%
  \pgfpathlineto{\pgfqpointpolar{198}{0.55\pgfplotmarksize}}%
  \pgfpathlineto{\pgfqpointpolar{234}{1.35\pgfplotmarksize}}%
  \pgfpathlineto{\pgfqpointpolar{270}{0.55\pgfplotmarksize}}%
  \pgfpathlineto{\pgfqpointpolar{306}{1.35\pgfplotmarksize}}%
  \pgfpathlineto{\pgfqpointpolar{342}{0.55\pgfplotmarksize}}%
  \pgfpathlineto{\pgfqpointpolar{378}{1.35\pgfplotmarksize}}%
  \pgfpathlineto{\pgfqpointpolar{414}{0.55\pgfplotmarksize}}%
  \pgfpathclose\pgfusepathqfillstroke}

% ---------------------------------------------------------------- diagrams
\tikzset{
  sb diagram/.style={font=\sbtext, line width=0.5pt, >={Stealth[length=3.2pt, width=2.6pt]}},
  sb box/.style={draw=sbDark, fill=white, rounded corners=1.5pt, inner sep=3pt, align=left, anchor=north west, font=\sbboxtext, execute at begin node={\hyphenpenalty=10000\exhyphenpenalty=10000\relax}},
  sb boxA/.style={sb box, draw=sbA, line width=0.8pt},
  sb boxB/.style={sb box, draw=sbB, line width=0.8pt},
  sb group/.style={fill=sbPale, draw=none, rounded corners=2.5pt},
  sb head/.style={font=\sbtext\bfseries, inner sep=0pt, outer sep=0pt, anchor=north west},
  sb arrow/.style={->, draw=sbDark, line width=0.6pt},
  sb note/.style={font=\sbtext, text=sbDark, inner sep=1.5pt, align=left},
  sb step/.style={circle, fill=sbInk, text=white, font=\sffamily\bfseries\fontsize{6}{6}\selectfont, inner sep=0pt, minimum size=8pt},
}
% the SQL that differs between the interpretations: a tinted box behind code
\newcommand{\hl}[2]{{\setlength{\fboxsep}{0.8pt}\colorbox{#1f}{\vphantom{Ay}\sbcode #2}}}
